\chapter{Quantum currents at every arity}
\label{chap:all-arity-current-actions}

A quantum symmetry must act on every finite collection of local
sources. Its identities relate different numbers of insertions.
The local triangle anomaly fixes one term in those identities.
Higher wheels also contribute at a positive heat scale. Their
variations contain collapsed wheels, which must cancel with
the Lie brackets of adjacent sources.

For the free theory on a holomorphic surface, every connected
matter graph is a path or a cycle. This restriction makes a
complete construction possible. Paths give the classical current
maps. Cycles give their scalar quantum corrections. A uniform
Gaussian estimate controls every partial collision, and the
remaining local curvature defines a central extension of the
source algebra.

\section{Compact sources and the meaning of the series}

Set $X=\mathbb C^2$, with $z_j=x_j+iy_j$, and fix
\[
 \Omega=dz_1\wedge dz_2,\qquad
 dV_4=dx_1\wedge dy_1\wedge dx_2\wedge dy_2,\qquad
 d\bar z_1d\bar z_2\Omega=4dV_4.
\]
Products of differential forms mean wedge products.
Let $\mathfrak g$ be an ordinary finite-dimensional complex Lie
algebra and let $\rho:\mathfrak g\to\operatorname{End}_{\mathbb C}(V)$
be a representation on an ordinary finite-dimensional complex
vector space. The compactly supported parameter algebra is
\[
 \mathfrak g_c=\Omega_c^{0,\bullet}(X,\mathfrak g),\qquad
 d=\bar\partial,\qquad
 [\alpha X_0,\eta Y_0]=(\alpha\wedge\eta)[X_0,Y_0].
\]
The degree of a parameter is its antiholomorphic form degree.

Write the full second matter field as
$\beta_{\mathrm{full}}=\beta\wedge\Omega$. Thus the coefficient
fields are
\[
 \gamma\in\Omega^{0,\bullet}(X)\otimes V,\qquad
 \beta\in\Omega^{0,\bullet}(X)\otimes V^\vee[1].
\]
The shift acts on the coefficient space, with
$E[k]^j=E^{j+k}$. Component coordinates of form degree $q$
have ghost numbers $-q$ for $\gamma$ and $1-q$ for $\beta$.
Holomorphic volume forms always stand on the right.

A finite collection of source insertions requires only a finite
family of compact smooth forms. For homogeneous
$\alpha_j\in\mathfrak g_c^{p_j}$, introduce independent formal
parameters and put
\[
 A=\sum_{j=1}^m t_j\alpha_j,\qquad |t_j|=1-p_j.
\]
The coefficient algebra $B$ is freely graded commutative on the
$t_j$. Odd parameters square to zero. The total degrees of
$\beta,A,\gamma$ are $1,1,0$, respectively, where total degree
adds ghost number to antiholomorphic form degree. The product is
\[
 (u\alpha)(v\eta)
   =(-1)^{|\alpha|\operatorname{gh}(v)}uv\,\alpha\wedge\eta.
\]
Odd coordinate derivatives act from the left.
The totalized spatial differential satisfies
$d(u\alpha)=(-1)^{\operatorname{gh}(u)}u\,d\alpha$ for formal
coefficients. The same sign applies to spatial differentiation
of coordinate-dependent components. Holomorphic derivatives
$\partial_i=\partial_{z_i}$ have degree zero.

On observables, the free differential $Q$ is the graded dual
of the field differential. In these conventions,
$Q\beta=d\beta$ and $Q\gamma=d\gamma$.
The Chevalley--Eilenberg differential on backgrounds is
\[
 d_{\mathrm{CE}}A=dA-A^2,\qquad
 d_{\mathrm{CE}}=d_{\mathrm{CE},1}+d_{\mathrm{CE},2}.
\]
The subscripts denote its linear and quadratic parts.
In a matrix product, $A$ denotes $\rho(A)$.
The source interaction
\[
 I(A)=\int_X\beta A\gamma\wedge\Omega
\]
has ghost degree zero and satisfies
\begin{equation}\label{eq:aac-bare-cme}
 (Q+d_{\mathrm{CE}})I+\frac12\{I,I\}_0=0,\qquad
 \frac12\{I,I\}_0=-\int_X\beta A^2\gamma\wedge\Omega.
\end{equation}
These are the local identities of
Chapter~\ref{chap:quantum-current-anomaly}. The bracket at zero
in this equation is the displayed local contraction, whose
diagonal evaluation exists.

The currents produced by these contractions need distribution
kernels. Use the observable space
\[
 \mathcal O_{\mathrm{dist}}
 =\bigoplus_{r\geq0}
 \mathcal E'_c\bigl(X^r;(E^\vee)^{\boxtimes r}
                     \otimes\operatorname{Dens}\bigr)_{S_r},
\]
Here $E$ is the graded bundle of matter components, and
$\operatorname{Dens}$ is the density bundle on $X^r$.
The summand $r=0$ is $\mathbb C$. A compact distribution is a
linear functional on smooth tests with an estimate
\[
 |T(\phi)|\leq C
       \max_{|\nu|\leq M}\sup_K|\partial^\nu\phi|
\]
for some compact $K$, finite $M$, and constant $C$.
Derivatives use real coordinates and a finite component basis.
The strong topology is uniform convergence on bounded families
of smooth tests. Such a family bounds every derivative on
each fixed compact set.

Multiplication uses external tensor products followed by graded
symmetrization. A smooth kernel contracts two selected slots by
\[
 (\mathcal C_{ij}T)(\psi)
       =T[P_{ij}\pi_{ij}^*\psi],
\]
where $\pi_{ij}$ forgets those slots and $P_{ij}$ includes their
pairing bivector, inserted with its graded permutation sign.
The finite-order estimate proves that the result is a compact
distribution. This construction, and the external-product
proof in Chapter~\ref{chap:quantum-current-anomaly}, apply to
joint kernels which need not split into linear measurements.

Fix a compact set $K\subset X$ containing the supports of the
chosen backgrounds. All continuity assertions below concern
bounded smooth families with this fixed support. On these
sectors, the displayed seminorm bounds control both background
coefficients and observable tests. No joint continuity on the
unrestricted global inductive-limit products is asserted.

Introduce an even variable $a$ of degree zero to record the
number of background insertions. The formal algebra is
\begin{equation}\label{eq:aac-formal-carrier}
 \bigl(\mathcal O_{\mathrm{dist}}\otimes B\bigr)
       [\hbar,\hbar^{-1}][[a]].
\end{equation}
Each coefficient of $a$ has finite matter degree and finite
Laurent support in the loop parameter $\hbar$.
The lower Laurent bound may depend on background order.
The topology means coefficientwise agreement below each fixed
power of $a$. There is no assertion of numerical convergence
of the background series.
In this notation the background differential is
\[
 d_{\mathrm{CE}}^{(a)}
       =d_{\mathrm{CE},1}+a\,d_{\mathrm{CE},2}.
\]
To evaluate an identity, one may enlarge the finite family by
the finitely many $d\alpha_i$ and $[\alpha_i,\alpha_j]$ that occur.
Its linear span need not be closed under either operation.

\section{Paths and cycles with their coefficients}

The regulated kernels are
\[
 \begin{split}
 N(u)&=\bar u_1d\bar u_2-\bar u_2d\bar u_1,\\
 p_t(u)&=\frac{e^{-|u|^2/(4t)}}{64\pi^2t^3}N(u),\\
 H_t(u)&=\frac{e^{-|u|^2/(4t)}}{64\pi^2t^2}
                           d\bar u_1d\bar u_2,\\
 P_{\epsilon,L}(u)&=\int_\epsilon^L p_t(u)\,dt,
                         \qquad 0<\epsilon<L.
 \end{split}
\]
They use full differences, including their differentials.
Thus $P_{\epsilon,L}(-u)=P_{\epsilon,L}(u)$ under pullback.
Their normalizations give
$dP_{\epsilon,L}=H_\epsilon-H_L$ and
$H_t\Omega=k_t\,dV_4$, where
$k_t=(4\pi t)^{-2}e^{-|u|^2/(4t)}$.

Write $\mathcal C_{\epsilon,L}=\partial_{P_{\epsilon,L}}$
for the even propagator contraction and
$\Delta_t=-\partial_{K_t}$ for the Batalin--Vilkovisky (BV)
heat contraction. The heat tensor $K_t$ includes the matter
pairing. On the algebra of joint distribution kernels,
\begin{equation}\label{eq:aac-contraction-identities}
 [Q,\mathcal C_{\epsilon,L}]=\Delta_\epsilon-\Delta_L,\qquad
 [\mathcal C_{\epsilon,L},\Delta_t]=0,\qquad
 Q\Delta_t+\Delta_tQ=\Delta_t^2=0.
\end{equation}
Transposition onto a smooth test proves the first equation.
Differentiation of the remaining test cancels, leaving
$-dP=-H_\epsilon+H_L$.
In two successive contractions, the selected pairs of slots
are disjoint. Their graded interchange proves the second
equation and the heat-square identity. The identity $dH_t=0$
proves compatibility with $Q$. These arguments use smooth
positive-cutoff kernels throughout.

The connected effective interaction is the formal expression
\begin{equation}\label{eq:aac-connected}
 I_{\epsilon,L}
 =\hbar\log\left(
    e^{\hbar\mathcal C_{\epsilon,L}}e^{aI/\hbar}\right).
\end{equation}
At background order $n$, the input has $2n$ matter legs, so
at most $n$ contractions act. The logarithm exists since its
argument equals $1+O(a)$.
Expanding a graph into its connected components proves the
exponential formula: the factorial for their unordered
collection exponentiates the sum of connected graphs.
Consequently the logarithm retains exactly those connected
graphs.

Every vertex has one $\beta$ leg and one $\gamma$ leg.
Following either leg therefore gives a directed path or a
directed cycle. A connected graph with $n$ vertices and $e$
edges has power $\hbar^{e-n+1}$.
Paths have two uncontracted matter legs and power zero.
Cycles have no matter legs and power one.

Set
\[
 s_n=(-1)^{(n-1)(n-2)/2},\qquad s_1=1.
\]
For $P=P_{\epsilon,L}$, define the ordered path functional
\begin{equation}
 \begin{split}
 T_n^{\epsilon,L}(A)
 =s_n\int_{X^n}
 &P(z_1-z_2)\cdots P(z_{n-1}-z_n)\\
 &{}\cdot\beta(z_1)A(z_1)\cdots A(z_n)\gamma(z_n)
                         \,\Omega_1\cdots\Omega_n.
 \end{split}
\end{equation}
The empty edge product for $n=1$ is one.
For $n\geq3$, define
\begin{equation}
 \begin{split}
 W_n^{\epsilon,L}(A)
 =\frac{s_n}{n}\int_{X^n}
 &P(z_1-z_2)\cdots P(z_n-z_1)\\
 &{}\cdot\operatorname{Tr}_V(A(z_1)\cdots A(z_n))
                         \,\Omega_1\cdots\Omega_n.
 \end{split}
\end{equation}
Put $W_1=W_2=0$. These are kernel integrals.
No infinite-dimensional operator trace is used.

The coefficients follow from left odd differentiation.
Start with $\prod_{j=1}^n(\beta_jA_j\gamma_j)$ and successively
pair $\gamma_i$ with $\beta_{i+1}$.
At step $i$, the left $\beta_{i+1}$ derivative crosses the
$i-1$ previous propagators, $\beta_1$, and $i$ backgrounds.
These are $2i$ odd factors, so their sign is positive.
The propagators accumulate in reverse order. Reversing the
$n-1$ odd propagators gives $s_n$.

There are $n!$ directed paths on $n$ labeled vertices.
Their $(n-1)!$ contraction orders cancel the denominator
$n!(n-1)!$ from the two exponentials.
Thus the path coefficient is $s_n$.
Closing the path crosses $n-1$ propagators with the odd
derivative. Moving the new propagator into edge order crosses
them again. The two signs cancel.
There are $(n-1)!$ directed cycles and $n!$ contraction orders.
Division by $n!n!$ leaves $s_n/n$.

Cyclic relabeling moves an odd propagator and an odd background
together, so it introduces no further sign.
The one-vertex loop is zero because $P(0)=0$.
The two-vertex loop contains $P(z-w)^2=0$.
Therefore
\begin{equation}
 I_{\epsilon,L}
   =\sum_{n\geq1}a^n
                  \bigl(T_n^{\epsilon,L}+\hbar W_n^{\epsilon,L}\bigr).
\end{equation}
There are no connected powers $\hbar^j$ with $j\geq2$
for these bilinear source vertices.

The same expressions can be written with interleaved kernels:
\[
 \begin{split}
 T_n&=\int\beta_1A_1P_{12}A_2\cdots
                      P_{n-1,n}A_n\gamma_n\,\Omega_1\cdots\Omega_n,\\
 W_n&=-\frac1n\int\operatorname{Tr}_V
       (A_1P_{12}A_2P_{23}\cdots A_nP_{n1})
                                      \,\Omega_1\cdots\Omega_n.
 \end{split}
\]
Moving the propagators left crosses
$\sum_{i=1}^{n-1}(i+1)$ odd factors for the path and
$\sum_{i=1}^n i$ for the cycle.
These parities reproduce $s_n$ and $s_n/n$, respectively.

\section{Every path has a compact observable kernel}

The edge differences
$u_i=z_i-z_{i+1}$ for $1\leq i<n$, together with $z_n$,
are independent coordinates. Their real Jacobian has absolute
value one. Each path propagator depends on one edge coordinate.
Use the Euclidean coefficient norm in the basis
$d\bar u_1,d\bar u_2$, so $|N(u)|=|u|$.
Polar coordinates give
\[
 \begin{split}
 \|p_t\|_{L^1}
 &=\frac{2\pi^2}{64\pi^2t^3}
        \int_0^\infty r^4e^{-r^2/(4t)}\,dr\\
 &=\frac{3\sqrt\pi}{8}t^{-1/2}.
 \end{split}
\]
The last integral follows by $r=2\sqrt t\,v$ and two
integrations by parts in the ordinary Gaussian integral.
Hence, for a constant $C$ independent of the cutoffs,
\begin{equation}\label{eq:aac-path-l1}
 \|P_{0,L}\|_{L^1}\leq C\sqrt L,\qquad
 \|P_{0,L}-P_{\epsilon,L}\|_{L^1}\leq C\sqrt\epsilon.
\end{equation}
Here $P_{0,L}=\int_0^L p_t\,dt$ has locally integrable
coefficients of order $|u|^{-3}$ near zero.

For $n\geq2$, telescope the difference of the $n-1$ edge
products, replacing one edge at a time.
Fubini's theorem and \eqref{eq:aac-path-l1} give
\begin{equation}
 \left\|
   \bigotimes_{i=1}^{n-1}P_{0,L}
   -\bigotimes_{i=1}^{n-1}P_{\epsilon,L}
 \right\|_{L^1}
 \leq C_n\sqrt\epsilon\,L^{(n-2)/2}.
\end{equation}
The bound is valid at all partial spatial collisions because
the edge coordinates are independent.

The background at $z_n$ restricts the base point to $K$.
Multiply by the other smooth background coefficients and push
forward the internal vertices. The result is a quadratic
matter kernel supported in $K\times K$.
The pushforward preserves the total-variation estimate.
Thus $T_n^{\epsilon,L}$ converges to a compact distribution
$T_n^L$ in total variation, uniformly on bounded background
families in this support sector.
Its size is bounded by
$C_{n,K}L^{(n-1)/2}\|\Phi\|_{C^0}$, where $\Phi$ includes
the backgrounds and the smooth test for the two remaining
matter slots. The difference obeys
\[
 |(T_n^L-T_n^{\epsilon,L})(\Phi)|
 \leq C_{n,K}\sqrt\epsilon\,L^{(n-2)/2}\|\Phi\|_{C^0}.
\]
A fixed compact cutoff may be inserted before taking this
seminorm. The term $T_1=I$ is independent of the cutoff.
Every tree limit has therefore been constructed from integrable
coefficients, without applying a singular kernel to an
arbitrary distribution.

\section{A Gaussian bound for all wheels}

For a cycle with $n\geq3$ edges, assign positive times
$t_1,\ldots,t_n$, and set
\[
 u_i=z_i-z_{i+1}\quad(i<n),\qquad
 u_n=-\sum_{i<n}u_i,\qquad T=\sum_{i=1}^n t_i.
\]
Use $u_1,\ldots,u_{n-1}$ as independent relative coordinates
and $z_n$ as the base point. Put
$D=\operatorname{diag}(t_1,\ldots,t_{n-1})$ and
$t=(t_1,\ldots,t_{n-1})^{\mathsf T}$.
For either complex coordinate, the Gaussian precision matrix is
\[
 M=\frac14\left(D^{-1}
                +t_n^{-1}\mathbf1\mathbf1^{\mathsf T}\right).
\]
The covariance is
\begin{equation}\label{eq:aac-covariance}
 C=M^{-1}=4\left(D-\frac{tt^{\mathsf T}}T\right).
\end{equation}
Indeed, $D^{-1}t=\mathbf1$ and $\mathbf1^{\mathsf T}t=T-t_n$.
Multiplying the two displayed matrices leaves the rank-one
coefficient
$-1/T+1/t_n-(T-t_n)/(t_nT)=0$.

Expand the determinant in its columns. Two columns taken from
the rank-one term are proportional, so their contribution is
zero. The remaining terms give
\[
 \det C
 =4^{n-1}\det D\left(1-\frac{\sum_{i<n}t_i}{T}\right)
 =4^{n-1}\frac{\prod_i t_i}{T}.
\]
There are two independent complex coordinate directions.
The relative Gaussian integral is consequently
\begin{equation}
 Z_n=(4\pi)^{2(n-1)}
                   \frac{(\prod_i t_i)^2}{T^2}.
\end{equation}
This follows by diagonalizing the positive matrix $M$ and
using $\int_{\mathbb C}e^{-\lambda|w|^2}dV_w=\pi/\lambda$
in each coordinate.

Let $\mathbb E_{\mathbf t}$ denote normalized Gaussian
expectation in the relative variables. Integration by parts
against a smooth compact coefficient gives, for $\mu=1,2$,
\[
 \begin{aligned}
 \mathbb E_{\mathbf t}[\bar u_{i,\mu}\Phi]
   &=\sum_{j<n}C_{ji}\,
                  \mathbb E_{\mathbf t}[\partial_{u_{j,\mu}}\Phi],
                                      &&i<n,\\
 \mathbb E_{\mathbf t}[\bar u_{n,\mu}\Phi]
   &=-\sum_{j<n}\frac{4t_nt_j}{T}\,
                  \mathbb E_{\mathbf t}[\partial_{u_{j,\mu}}\Phi].
 \end{aligned}
\]
For the second formula, substitute $u_n=-\sum_{i<n}u_i$
and sum the columns of \eqref{eq:aac-covariance}.
Their absolute sums satisfy
\begin{equation}\label{eq:aac-covariance-bounds}
 \sum_{j<n}|C_{ji}|\leq8t_i,\qquad
 \sum_{j<n}\frac{4t_nt_j}{T}\leq4t_n.
\end{equation}
For the first inequality, bound the diagonal contribution by
$4t_i$ and the sum of the off-diagonal contributions by
$4t_i\sum_{j<n}t_j/T$.
No ratio between different times enters these bounds.

Expand $N(u_1)\cdots N(u_n)$ in a relative exterior basis.
Every coefficient contains one antiholomorphic linear coordinate
from each edge. Apply the integration-by-parts formulas to all
$n$ factors. Holomorphic derivatives annihilate the remaining
antiholomorphic numerators. All derivatives therefore fall on
the smooth external coefficient.
Equation~\eqref{eq:aac-covariance-bounds} supplies
$C_n\prod_i t_i$ times its $C^n$ seminorm.

The absolute scalar prefactor, including the relative volume
conversion, is
\[
 \frac{4^{n-1}Z_n}{(64\pi^2)^n\prod_i t_i^3}
       =\frac1{64\pi^2}\frac1{T^2\prod_i t_i}.
\]
The numerator factors cancel every individual $t_i$ denominator.
For each component of the wheel distribution this proves
\begin{equation}\label{eq:aac-wheel-estimate}
 |\mathcal W_{\mathbf t}(\Phi)|
       \leq \frac{C_{n,K}}{(t_1+\cdots+t_n)^2}
                                    \|\Phi\|_{C^n}.
\end{equation}
The constant includes the finite exterior expansion, color
contraction, and volume of the compact base projection.
It may depend on $n$.

The estimate holds throughout the positive time orthant.
It controls faces where only some times vanish, as well as
the common time origin. The slice of the positive simplex
$\sum_i t_i=\tau$ has measure $\tau^{n-1}/(n-1)!$.
Since the time cube lies in $\sum_i t_i\leq nL$,
\begin{equation}\label{eq:aac-wheel-time-bound}
 \int_{[0,L]^n}\frac{d^nt}{(\sum_i t_i)^2}
 \leq\frac1{(n-1)!}\int_0^{nL}\tau^{n-3}\,d\tau
 =\frac{(nL)^{n-2}}{(n-1)!(n-2)}.
\end{equation}
This proves the lower-cutoff limit of every wheel.

The excluded region is $\{\min_i t_i<\epsilon\}$.
It lies in the union of $n$ slabs.
For $n\geq4$, integrate the other $n-1$ times first and
apply \eqref{eq:aac-wheel-time-bound} in that dimension.
The error is bounded by $C_n\epsilon L^{n-3}$.
For $n=3$, the integral over the other two times is
\[
 \int_0^L\!\int_0^L\frac{ds\,dt}{(u+s+t)^2}
       =\log\frac{(u+L)^2}{u(u+2L)}.
\]
Its integral over $0<u<\epsilon$ is bounded by
$C\epsilon(1+\log(L/\epsilon))$.
Together with \eqref{eq:aac-wheel-estimate}, these estimates
give strong convergence on bounded compact test families.
The limiting wheel obeys
\[
 |W_n^L(\Phi)|\leq C_{n,K}L^{n-2}\|\Phi\|_{C^n}.
\]

The universal kernel is a distribution on $X^n$, generally
with noncompact support because of translation invariance.
Its scalar components belong to $\mathcal D'(X^n)$, the
linear functionals on compactly supported smooth tests with
a finite-order estimate on each fixed compact support.
Multiplication by fixed compact background cutoffs gives
compact distributions. Evaluation on the chosen backgrounds
gives a scalar matter observable $W_n^L(A)$.
The proof constructs the regulated distributional limit.
It does not require absolute integrability of the unregularized
product of singular propagators.

Applying $d_{\mathrm{CE},1}$ costs one more spatial derivative
in the seminorm. A derivative in a finite background parameter
replaces an insertion, so it gives finitely many terms with
the same majorant. The quadratic differential multiplies or
brackets two smooth insertions. The finite Leibniz expansion
controls all derivatives of that product.
Thus every differential used below preserves these
support-sector bounds.

\section{Heat edges and collapsed insertions}

The local source identity uses a diagonal contraction, whereas
the effective bracket uses a heat kernel. Their difference
places one distinguished edge $H_\epsilon-H_0$ in a graph.
Here $H_0\Omega=\delta_0dV_4$ is the normalized diagonal current.
For paths and sufficiently long cycles, the two parts have
the same limit.

In a marked path, choose the marked difference $v$ as one
independent edge coordinate. After its top form has selected
the external component, the heat-minus-diagonal term is
\[
 \int k_\epsilon(v)
             \bigl(\Phi(v,u)-\Phi(0,u)\bigr)\,dV_v.
\]
The Gaussian first moment is at most $C\sqrt\epsilon$.
The mean-value theorem and the $n-2$ remaining path estimates
give
\begin{equation}\label{eq:aac-marked-path}
 |\mathcal M_{\mathrm{path},\epsilon}(\Phi)
                 -\mathcal M_{\mathrm{path},0}(\Phi)|
 \leq C_{n,K}\sqrt\epsilon\,
                  L^{(n-2)/2}\|\Phi\|_{C^1}.
\end{equation}
This proves strong convergence to zero of every marked path
difference.

For a cycle, put the heat edge in the closing position, with
time $r>0$. The other $n-1$ edges carry propagators.
Its heat denominator is $r^2$, and there are only $n-1$
antiholomorphic numerator factors.
Repeat the covariance calculation with $t_n=r$.
Before integration by parts, the scalar prefactor is
\[
 \frac1{64\pi^2}
 \frac1{(r+\sum_{i<n}t_i)^2\prod_{i<n}t_i}.
\]
The $n-1$ numerator integrations by parts cancel the product
in the denominator. Hence
\begin{equation}\label{eq:aac-marked-wheel}
 |\mathcal M_{r;\mathbf t}(\Phi)|
 \leq \frac{C_{n,K}}{(r+\sum_{i<n}t_i)^2}
                                      \|\Phi\|_{C^{n-1}}.
\end{equation}
For $n\geq4$, this bound is integrable uniformly as $r\downarrow0$:
\[
 \int_{[0,L]^{n-1}}\frac{d^{n-1}t}{(\sum_{i<n}t_i)^2}
 \leq\frac{((n-1)L)^{n-3}}{(n-2)!(n-3)}.
\]

At fixed positive propagator times, the heat approximate
identity collapses the marked edge with coefficient $+1$.
The adjacent backgrounds are multiplied in their existing
order, with the original labels and graded signs retained.
The result is an unmarked $(n-1)$-cycle with one composite
insertion.

Strong convergence requires uniformity on the test family.
Split the time cube into $[\delta,L]^{n-1}$ and its complement.
On the first region, the propagators and all required
derivatives are uniformly bounded, so the heat error tends
uniformly to zero. On the complement, the integrable bound
in \eqref{eq:aac-marked-wheel} is uniformly small as
$\delta\downarrow0$.
The collapsed term has the same bound by its fixed-time limit.
Consequently
\begin{equation}\label{eq:aac-collapsed-difference}
 \mathcal M_{\epsilon;[\epsilon,L]^{n-1}}
       -\mathcal M_{0;[\epsilon,L]^{n-1}}
       \longrightarrow0,\qquad n\geq4.
\end{equation}
This is the difference of the two marked terms.
The individual heat-marked limit can be nonzero and nonlocal
at fixed $L$.

For a four-cycle, collapse edge $41$ and set
$x=z_1-z_3$ and $y=z_2-z_3$.
The remaining numerator is $N(x-y)N(y)N(x)$.
Its coefficient of $d\bar x_1d\bar x_2d\bar y_1$ equals
\[
 \bar y_2(\bar x_1\bar y_2-\bar x_2\bar y_1).
\]
It equals one at $x=(1,0)$ and $y=(0,1)$.
The scalar heat-time factors there are positive.
A complementary $d\bar y_2$ component and a base $(0,2)$-form
therefore give a nonzero smooth coefficient off every remaining
edge diagonal. Compact bumps in sufficiently small neighborhoods
give a nonzero value. Take $\mathfrak g=\mathfrak{gl}_4$ and
$V=\mathbb C^4$ with its defining representation.
The ordered color matrices $E_{12},E_{23},E_{34},E_{41}$
have trace product one.
They isolate the corresponding directed color cycle.
Thus vanishing of an individual higher heat-marked wheel
would contradict an explicit compact evaluation.

The one-vertex heat contraction vanishes under diagonal
pullback of the full difference form.
The two-vertex marked cycle has relative antiholomorphic
degree three in two relative complex coordinates, so it
vanishes.
The three-vertex cycle is exceptional because
$(s+t)^{-2}$ is not integrable at the origin of its two
propagator times. Its exterior determinant supplies a local
term instead of the collapsed limit.

\section{The exact finite-scale Ward identity}

Define the bare mismatch
\[
 B_\epsilon(A)=\frac12
                  \bigl(\{I,I\}_\epsilon-\{I,I\}_0\bigr).
\]
The first bracket is a smooth heat contraction, and the second
is the explicit local bracket in \eqref{eq:aac-bare-cme}.
The heat tadpole $\Delta_\epsilon I$ is zero.
For
$D_\sigma=Q+d_{\mathrm{CE}}^{(a)}+\hbar\Delta_\sigma$,
the second-order product rule therefore gives
\[
 D_\epsilon e^{aI/\hbar}
        =\frac{a^2}{\hbar}B_\epsilon e^{aI/\hbar}.
\]
The background differential commutes with the propagator
contraction. Equation~\eqref{eq:aac-contraction-identities}
and commutation of the contraction operators give
\[
 D_L e^{\hbar\mathcal C_{\epsilon,L}}
       =e^{\hbar\mathcal C_{\epsilon,L}}D_\epsilon.
\]
To verify this directly, expand the exponential.
For $C=\mathcal C_{\epsilon,L}$, its commutator is
$[Q,C^j]=jC^{j-1}(\Delta_\epsilon-\Delta_L)$.
The resulting finite coefficient sums give the displayed
identity.

Let the quantum master equation (QME) defect be
\[
 \mathcal R_{\epsilon,L}
 =(Q+d_{\mathrm{CE}}^{(a)})I_{\epsilon,L}
   +\frac12\{I_{\epsilon,L},I_{\epsilon,L}\}_L
   +\hbar\Delta_LI_{\epsilon,L}.
\]
Apply $D_L$ to \eqref{eq:aac-connected} in exponential form
and divide by that exponential. The exact identity is
\begin{equation}\label{eq:aac-regulated-defect}
 \mathcal R_{\epsilon,L}
 =a^2e^{-I_{\epsilon,L}/\hbar}
       e^{\hbar\mathcal C_{\epsilon,L}}
                  \bigl(B_\epsilon e^{aI/\hbar}\bigr).
\end{equation}
No assertion that the positive-cutoff bare mismatch vanishes
is needed.

Normalization by $e^{-I_{\epsilon,L}/\hbar}$ removes connected
components disjoint from the distinguished insertion.
The remaining component is a path or cycle containing the
edge $H_\epsilon-H_0$.
Its coefficients can be fixed before any edge is collapsed.
Keep all $n$ original vertex labels. Symmetrizing the
distinguished two-vertex insertion gives
\begin{equation}\label{eq:aac-marked-count}
 \frac1{2(n-2)!}\Delta_{12}(I_1I_2)
             \prod_{j=3}^n I_j
 \quad\longmapsto\quad
 \frac1{n!}\sum_{i<j}\Delta_{ij}\prod_{k=1}^n I_k.
\end{equation}
Here $\Delta_{ij}$ denotes the distinguished cross-contraction,
with either endpoint kernel and the BV minus sign.
It includes both directed matter pairings.
The equality after symmetrization follows from
$\binom n2/n!=1/[2(n-2)!]$.

There are $n(n-1)!=n!$ labeled directed cycles with a chosen
marked edge. The remaining $(n-1)!$ contraction orders cancel
their exponential denominator. The marked-cycle coefficient
therefore has magnitude one.
To find its sign, close the ordered path with
$-H_{n1}\partial_{\beta_1}\partial_{\gamma_n}$ and place
$H_{n1}$ first. The derivative crosses $n-1$ odd propagators,
and the heat form has even degree. Thus the sign is
\[
 m_n=-(-1)^{n-1}s_n
       =(-1)^{n(n-1)/2+1}.
\]
In particular, $m_3=+1$, while $m_4=m_5=-1$.

The $H_\epsilon$ and $H_0$ terms in
\eqref{eq:aac-marked-count} have identical labels and
derivative orders. Diagonal collapse changes neither their
coefficient nor their inherited sign.
Its local ordered insertion is exactly
$-\int\beta A^2\gamma\wedge\Omega$.
Equations~\eqref{eq:aac-marked-path} and
\eqref{eq:aac-collapsed-difference} therefore cancel every
marked path difference and every marked cycle difference of
arity at least four.

For the triangle, use the relative coordinates and orientation
of Lemma~\ref{lem:qca-localization}. Its marked geometric
distribution is
\[
 -\frac1{4\pi^2}\int_\epsilon^L\!\int_\epsilon^L
    \frac{\epsilon\,ds\,dt}{(\epsilon+s+t)^3}
    \int_X\mathbb E[
       (\partial_{x_1}\partial_{y_2}
        -\partial_{x_2}\partial_{y_1})\Phi]\,
                                  d\bar z_1d\bar z_2\Omega_z.
\]
The exterior numerator is
$\bar x_2\bar y_1-\bar x_1\bar y_2$.
Its Gaussian integration-by-parts factor is
$-16\epsilon st/(\epsilon+s+t)$.
Together with the Gaussian mass and relative volume factor
$16$, this gives the displayed coefficient.
The time mass equals
\[
 \int_\epsilon^L\!\int_\epsilon^L
 \frac{\epsilon\,ds\,dt}{(\epsilon+s+t)^3}
 =\frac16-\frac{\epsilon}{L+2\epsilon}
                  +\frac{\epsilon}{2(2L+\epsilon)}.
\]
It tends to $1/6$ at every fixed $L>0$.
The covariance first moment gives localization error
$C_\Phi\sqrt\epsilon+C'_\Phi\epsilon/L$.
The collapsed two-wheel is zero because it contains $P(u)^2$.

With the marked sign $m_3=+1$ and the color trace, the
surviving scalar polynomial is
\begin{equation}\label{eq:aac-theta}
 \Theta_\rho(A)=-\frac1{24\pi^2}
  \int_X\operatorname{Tr}_V
     \left[A(\partial_1A\,\partial_2A
                         -\partial_2A\,\partial_1A)\right]\wedge\Omega.
\end{equation}
This is a fixed-$L$ lower-cutoff limit.
The scalar triangle has no remaining matter leg to which
another source vertex could attach.

The left side of \eqref{eq:aac-regulated-defect} also converges
in the specified topology on the distribution kernels.
Transposition onto differentiated smooth tests makes $Q$
continuous in each support sector.
For a fixed positive-scale contraction and a bounded output
test family $B_0$, the inserted tests
$\{H_L\pi^*\psi:\psi\in B_0\}$ form a bounded family on
every required compact input support.
Thus the contraction commutes with the strong limits.

Products of the convergent tree kernels also converge.
Their supports are fixed and their total variations are
uniformly bounded. Apply
\[
 S_\epsilon\boxtimes T_\epsilon-S\boxtimes T
   =(S_\epsilon-S)\boxtimes T_\epsilon
                  +S\boxtimes(T_\epsilon-T)
\]
and the total-variation estimates to each term.
The same argument with the finite-order bounds controls the
bounded test families after differentiation.
Evaluated wheels are scalar in matter, so their brackets and
matter BV Laplacians vanish.

\begin{theorem}\label{thm:aac-fixed-scale}
For every $L>0$, the series
\[
 I_L=\sum_{n\geq1}a^n(T_n^L+\hbar W_n^L)
\]
is defined in \eqref{eq:aac-formal-carrier} on every finite
family of compact smooth backgrounds. It satisfies
\begin{equation}\label{eq:aac-all-order-qme}
 (Q+d_{\mathrm{CE}}^{(a)})I_L
       +\frac12\{I_L,I_L\}_L+\hbar\Delta_LI_L
       =\hbar a^3\Theta_\rho(A).
\end{equation}
The limits and operations obey the stated bounded-family
estimates in each fixed compact support sector.
\end{theorem}

\begin{proof}
The path and wheel estimates construct each coefficient.
The continuity arguments pass every term on the left of
\eqref{eq:aac-regulated-defect} to the limit.
The labeled marked-graph calculation cancels all terms on
its right except the scalar triangle.
Its limit is \eqref{eq:aac-theta}.

At background order $N$, every exponential or logarithmic
coefficient is a finite sum. The bracket uses only the
finitely many pairs $i+j=N$, and the quadratic background
differential uses order $N-1$.
Laurent support in $\hbar$ is finite at each such order.
Thus the coefficientwise limits require no interchange with
an infinite numerical series.
\end{proof}

Separate the quadratic and scalar matter terms, with
$T_0=W_0=0$. The theorem gives all identities
\begin{equation}\label{eq:aac-tree-identities}
 (Q+d_{\mathrm{CE},1})T_n^L
       +d_{\mathrm{CE},2}T_{n-1}^L
       +\frac12\sum_{\substack{i+j=n\\i,j\geq1}}
                             \{T_i^L,T_j^L\}_L=0,
\end{equation}
and
\begin{equation}
 d_{\mathrm{CE},1}W_n^L+d_{\mathrm{CE},2}W_{n-1}^L
          +\Delta_LT_n^L
  =\begin{cases}
      \Theta_\rho,&n=3,\\
      0,&n\ne3.
    \end{cases}
\end{equation}
The higher scalar corrections are part of the finite-scale
quantum symmetry, even though its local curvature is cubic.

\section{Transport between positive heat scales}

For $0<L_1<L_2$, the interval propagator $P_{L_1,L_2}$ is
smooth. It therefore acts on the constructed distribution
observables. The interactions satisfy
\begin{equation}\label{eq:aac-scale-transport}
 I_{L_2}
 =\hbar\log\left(
       e^{\hbar\mathcal C_{L_1,L_2}}e^{I_{L_1}/\hbar}\right).
\end{equation}
Keep a lower cutoff $\delta<L_1$ first.
The identity
$P_{\delta,L_2}=P_{\delta,L_1}+P_{L_1,L_2}$
and commutation of the even contractions prove
\eqref{eq:aac-scale-transport} before taking $\delta\to0$.
The path and wheel bounds, followed by smooth contraction,
justify that limit at each background order.

The curvature is scalar in matter. Smooth propagator
contraction fixes it, so transport preserves
$\hbar\Theta_\rho$.
Smooth dependence on $L$ follows on each compact interval of
positive scales. Choose $L_*>0$ below that interval and use
\eqref{eq:aac-scale-transport} from $L_*$.
All derivatives of $P_{L_*,L}$ are bounded on the required
compact supports. At fixed background order the contractions
and logarithm have finitely many terms.
Every scale derivative therefore passes through the
construction in the stated support sectors.
No derivative at the singular lower endpoint is needed.

\section{Suspension and the source differential}

The identities \eqref{eq:aac-tree-identities} have a multilinear
meaning: the current map preserves the parameter bracket up
to specified higher homotopies. Their signs depend on the
suspension convention. Fix that convention before extracting
the maps.

Let $\mathcal U=\mathfrak g_c[1]$, and write
$s:\mathfrak g_c\to\mathcal U$ for the degree-minus-one shift.
For $\alpha_i$ of degree $p_i$, put
$x_i=s\alpha_i$ and $r_i=|x_i|=p_i-1$.
All coalgebras below are over $\mathbb C$.
For a graded complex vector space $W$, put
\[
 S^c(W)=\bigoplus_{n\geq0}S^n(W),
 \qquad xy=(-1)^{|x||y|}yx.
\]
Its product is graded symmetric multiplication.
The empty word $1$ is the unit, the coaugmentation sends
$\lambda$ to $\lambda1$, and the counit
$\varepsilon:S^c(W)\to\mathbb C$ projects to length zero.
An $(i_1,\ldots,i_k)$-unshuffle is a permutation preserving the
order inside each successive block of those lengths.
Write $\operatorname{Sh}(i_1,\ldots,i_k)$ for this set.
For homogeneous $w_i\in W$, let $\epsilon_s(\sigma;w)$
be the Koszul sign obtained by swapping adjacent $w_i,w_j$
with factor $(-1)^{|w_i||w_j|}$.
When $w_i=x_i$, abbreviate this sign to $\epsilon_s(\sigma)$.
The unshuffle coproduct is
\[
 \Delta(w_1\cdots w_n)
 =\sum_{i=0}^n\ \sum_{\sigma\in\operatorname{Sh}(i,n-i)}
   \epsilon_s(\sigma;w)\,
   (w_{\sigma(1)}\cdots w_{\sigma(i)})
       \otimes(w_{\sigma(i+1)}\cdots w_{\sigma(n)}).
\]
In particular, $\Delta1=1\otimes1$.
Splitting a word into three ordered subsequences describes
both iterated coproducts with the same Koszul sign, proving
coassociativity. The two empty-block terms give the counit
identities.

On $\overline S^c(W)=\ker\varepsilon$, the reduced coproduct is
$\overline\Delta z=\Delta z-1\otimes z-z\otimes1$.
Its $n$fold iteration into $n+1$ nonempty blocks vanishes on
every word of length at most $n$.
This finite-length property is the conilpotence used here.
All the following statements concern finite words; formal
coefficients are treated coefficientwise.

\begin{lemma}\label{lem:aac-symmetric-coderivations}
The primitive elements of $S^c(W)$ are exactly $W$.
Let $f:S^c(W)\to S^c(Z)$ be a degree-zero coalgebra map
preserving the coaugmentation and counit.
A homogeneous coderivation along $f$ is a linear map $D$
satisfying $D1=0$ and
\[
 \Delta D=(D\otimes f+f\otimes D)\Delta.
\]
Here graded tensor products of maps mean
$(A\otimes B)(v\otimes z)=(-1)^{|B||v|}Av\otimes Bz$.
Such a coderivation is determined by its corestriction
$\pi_ZD$, where $\pi_Z:S^c(Z)\to Z$ projects to length one.
\end{lemma}

\begin{proof}
An element $z$ is primitive if
$\Delta z=z\otimes1+1\otimes z$.
For $m\geq2$, denote by $\Delta_{1,m-1}$ the component
$S^m(W)\to W\otimes S^{m-1}(W)$ of the coproduct, and by
$\mu$ graded symmetric multiplication. Then
\[
 \mu\Delta_{1,m-1}(w_1\cdots w_m)=m\,w_1\cdots w_m.
\]
There are $m$ choices of the first output.
Each summand's unshuffle sign is repeated when multiplication
reorders its factors to $w_1\cdots w_m$.
The two signs cancel, so every summand is the original word.
Since $m$ is invertible in $\mathbb C$, $\Delta_{1,m-1}$ is
injective. This identity also holds when the word vanishes
by graded symmetry.

Decompose a primitive into its finitely many length components.
Projection to bidegree $(1,m-1)$ kills its component of
length $m\geq2$ by that injectivity.
Its scalar component $\lambda1$ must also vanish:
$\Delta(\lambda1)=\lambda1\otimes1$, whereas primitivity
requires $2\lambda1\otimes1$.
Every element of $W$ is primitive by the coproduct formula.
This proves the first assertion.

For the second, subtract two coderivations with the same
corestriction. It suffices to suppose $\pi_ZD=0$ and prove
$D=0$. The empty word is killed by hypothesis.
Suppose $D$ kills words of length less than $n$.
For a word $v$ of length $n$, every nonempty proper split
in $\Delta v$ has both lengths less than $n$.
The coderivation identity and the induction hypothesis give
\[
 \Delta D(v)=D(v)\otimes1+1\otimes D(v).
\]
Thus $D(v)$ is primitive. It lies in $Z$, where
$\pi_ZD(v)=0$ forces $D(v)=0$.
Induction proves the assertion on every finite word.
\end{proof}

For completeness, homogeneous maps
$a_j:S^j(W)\to W$ of a common degree $h$, with $j\geq1$,
extend to a coderivation along the identity by
\[
 D_a(w_1\cdots w_n)
 =\sum_{j=1}^n\ \sum_{\sigma\in\operatorname{Sh}(j,n-j)}
   \epsilon_s(\sigma;w)\,
   a_j(w_{\sigma(1)}\cdots w_{\sigma(j)})
          w_{\sigma(j+1)}\cdots w_{\sigma(n)},
 \qquad D_a1=0.
\]
To check its coproduct identity, choose the block on which
$a_j$ acts and then split the remaining inputs into left
and right blocks. Its output must lie on one of those sides.
An output on the left gives exactly
$(D_a\otimes1)\Delta$.
For an output on the right, let $b$ be the degree of the
chosen input block and $\ell$ the degree of the left block.
Moving the output past the left block gives
$(-1)^{(b+h)\ell}$; moving the chosen inputs past that block
in the other order gives $(-1)^{b\ell}$.
The quotient is $(-1)^{h\ell}$, exactly the tensor sign in
$(1\otimes D_a)\Delta$.
These choices exhaust both sides with equal multiplicities.
Hence the displayed extension is a coderivation, and its
corestrictions are the $a_j$.

Two consequences will be used repeatedly.
An odd coderivation $q$ along the identity has a square
which is again a coderivation, since
\[
 (q\otimes1+1\otimes q)^2=q^2\otimes1+1\otimes q^2.
\]
The mixed products are $q\otimes q$ and $-q\otimes q$.
Also, if $q,q'$ are coderivations on source and target,
then $q'f$ and $fq$ are coderivations along $f$:
substitute $\Delta f=(f\otimes f)\Delta$ into their
coproducts. Their difference is therefore a coderivation
along $f$ as well.
The lemma proves that a zero corestriction kills each of
these maps in full.

Specify the degree-one coderivation $q$ by its components
to $\mathcal U$:
\begin{equation}\label{eq:aac-source-q}
 q_1(s\alpha)=s(d\alpha),\qquad
 q_2(s\alpha\,s\eta)=(-1)^p s[\alpha,\eta],
       \qquad |\alpha|=p.
\end{equation}
There are no higher source components yet.
A component $q_j$ acts on a word by summing over
$(j,n-j)$-unshuffles, applying $q_j$ to the first block and
multiplying its output by the remaining block.
This defines the full coderivation.

The binary formula is graded symmetric on $\mathcal U$:
interchanging $\alpha,\eta$ combines graded antisymmetry of
their Lie bracket with the shift sign.
The unary square is zero because $d^2=0$.
At arity two, the three terms are
\[
 \begin{split}
 q_1q_2(s\alpha,s\eta)
    &=(-1)^p s[d\alpha,\eta]+s[\alpha,d\eta],\\
 q_2(q_1s\alpha,s\eta)
    &=(-1)^{p+1}s[d\alpha,\eta],\\
 (-1)^{p-1}q_2(s\alpha,q_1s\eta)
    &=-s[\alpha,d\eta].
 \end{split}
\]
They cancel. At arity three, the corestriction of $q^2$ is
the graded Jacobi sum, which vanishes.
These exhaust its possible corestrictions.
The square is a coderivation, so
Lemma~\ref{lem:aac-symmetric-coderivations} gives $q^2=0$.
In particular, the chosen unary operation is $+sd$.
It is not the canonical differential $-sd$ on the shifted
complex considered in isolation.

Chevalley--Eilenberg cochains are graded multilinear functions
on these finite tensors, completed in background weight when
needed. Their differential and product are
\begin{equation}\label{eq:aac-dual-conventions}
 \begin{split}
 (d_{\mathrm{CE}}\phi)(v)
    &=-(-1)^{|\phi|}\phi(qv),\\
 (\phi\psi)(v)
    &=\sum_{\Delta v}
       (-1)^{|\psi||v_{(1)}|}
                    \phi(v_{(1)})\psi(v_{(2)}).
 \end{split}
\end{equation}
The coproduct sign is included in $\Delta v$.
These formulas apply to labeled input positions even when
their images in $\mathfrak g_c$ are linearly dependent.

They induce the background differential used above.
For the linear term, the coordinate corresponding to
$d\alpha$ has degree $-p$. Thus
\[
 (d_{\mathrm{CE}}t_{d\alpha})(s\alpha)
       =(-1)^{p+1},
\]
which is the coefficient of the totalized $dA$.
For the quadratic term, ghost-first multiplication gives
\[
 [t_\alpha t_\eta](-A^2)
           =(-1)^{1+p+pq}[\alpha,\eta].
\]
Evaluating the two coordinate factors on
$s\alpha\,s\eta$ adds $(-1)^{(p-1)(q-1)}$.
The result is $(-1)^q[\alpha,\eta]$.
Dualizing the second formula in \eqref{eq:aac-source-q}
gives the same sign, since the bracket coordinate has degree
$1-p-q$. Therefore these conventions give
$d_{\mathrm{CE}}A=dA-A^2$ in both its linear and quadratic
parts.

For the quantum target, put
\[
 \mathcal L_L=\mathcal O_{\mathrm{dist}}[[\hbar]][-1],
 \qquad |uF|=|F|+1,\qquad D_L=Q+\hbar\Delta_L.
\]
Its differential and bracket are
\[
 d_{\mathcal L_L}(uF)=-uD_LF,\qquad
 [uF,uG]_L=u\{F,G\}_L.
\]
The BV bracket has degree one on observables, so this is a
degree-zero Lie bracket on $\mathcal L_L$.
Identify the suspension by $s'(uF)=F$.
The target coderivation has components
\begin{equation}\label{eq:aac-target-q}
 q'_1(F)=-D_LF,\qquad
 q'_2(F,G)=(-1)^{|F|+1}\{F,G\}_L.
\end{equation}
The equations $D_L^2=0$, the odd-bracket differential
identity, and the odd Jacobi identity kill the corestrictions
of $(q')^2$ at arities one, two, and three.
There are no others. Its square is a coderivation, so
Lemma~\ref{lem:aac-symmetric-coderivations} gives $(q')^2=0$.
For the classical target, replace $D_L$ by $Q$.
All these operations use positive-scale contractions.

\section{The Taylor components in every degree}

For homogeneous compact inputs define the ordered kernels
\[
 \begin{split}
 K_{n,L}(\alpha_1,\ldots,\alpha_n)
 =\int_{X^n}
 &P_{0,L,12}\cdots P_{0,L,n-1,n}\\
 &{}\cdot\beta_1\rho(\alpha_1)_1\cdots
             \rho(\alpha_n)_n\gamma_n\,\Omega_1\cdots\Omega_n,
 \end{split}
\]
and, for $n\geq3$,
\[
 \begin{split}
 C_{n,L}(\alpha_1,\ldots,\alpha_n)
 =\lim_{\epsilon\to0}\int_{X^n}
 &P_{\epsilon,L,12}\cdots P_{\epsilon,L,n1}\\
 &{}\cdot\operatorname{Tr}_V
     \bigl(\rho(\alpha_1)_1\cdots\rho(\alpha_n)_n\bigr)
                                      \,\Omega_1\cdots\Omega_n.
 \end{split}
\]
Set $C_{1,L}=C_{2,L}=0$.
The preceding estimates define these integrals and limits.
The first is a compact quadratic matter kernel after
background evaluation. The second is scalar in matter.

For degrees $p_i=|\alpha_i|$, let $\epsilon(\sigma;p)$ be
the Koszul sign for permuting the unshifted inputs.
Define the unnormalized graded antisymmetrization by
\[
 \operatorname{Alt}K(\alpha_1,\ldots,\alpha_n)
 =\sum_{\sigma\in S_n}\operatorname{sgn}(\sigma)
      \epsilon(\sigma;p)\,
        K(\alpha_{\sigma(1)},\ldots,\alpha_{\sigma(n)}).
\]
There is no $1/n!$ in this definition.
The classical and quantum component formulas are
\begin{equation}\label{eq:aac-current-components}
 \begin{split}
 F_n^{\mathrm{cl}}(\alpha_1,\ldots,\alpha_n)
    &=(-1)^{\sum_i p_i+1}\,
                         u\,\operatorname{Alt}K_{n,L},\\
 F_n^q(\alpha_1,\ldots,\alpha_n)
    &=(-1)^{\sum_i p_i+1}\,
       u\left(\operatorname{Alt}K_{n,L}
                      +\frac{\hbar}{n}\operatorname{Alt}C_{n,L}\right).
 \end{split}
\end{equation}
These are algebraic multilinear maps on compact smooth inputs.
Their support-sector continuity is exactly that of the kernel
estimates above.

The tree kernel has observable degree $\sum_i p_i-n$.
Indeed, integration selects
\[
 q_\beta+q_\gamma+\sum_i p_i+(n-1)=2n,
\]
and the matter ghost degree is $1-q_\beta-q_\gamma$.
A scalar wheel can survive only when $\sum_i p_i=n$.
The shift $u$ therefore makes both parts of $F_n^q$ have
degree $1-n$. The target contains only nonnegative powers
of $\hbar$, even though the exponential construction used
finite Laurent polynomials.

The precise passage from universal polynomials to these maps
accounts for both parameter ordering and graded duality.
Put
\[
 E_n(p)=\sum_{i<j}(p_i-1)(p_j-1),\qquad
 B_n(p)=\sum_i(n-i)p_i,\qquad
 \delta_n(p)=B_n(p)+E_n(p).
\]
For $A=\sum_{i=1}^n t_i\alpha_i$, the ordered coefficient
prescription is
\begin{equation}
 F_n^q(\alpha_1,\ldots,\alpha_n)
 =(-1)^{\delta_n(p)}u\,
       [t_1\cdots t_n](T_n^L+\hbar W_n^L)(A).
\end{equation}
The coefficient has the parameters in exactly the indicated
order. The product convention \eqref{eq:aac-dual-conventions}
gives
\[
 (t_1\cdots t_n)(x_1\cdots x_n)=(-1)^{E_n(p)}
\]
on the corresponding labeled inputs.
This is the second sign in $\delta_n$.

For an ordered tree, moving the parameters left crosses the
$n-1$ propagators and $\beta$, followed by the preceding
background forms. Its exponent is
\[
 M_n(p)=n\sum_i(1-p_i)+\sum_{i<j}p_i(1-p_j).
\]
For a wheel, $n$ propagators replace those $n$ odd factors,
so the exponent is unchanged.
Let $S_n=(n-1)(n-2)/2$ and $P_n=\sum_{i<j}p_ip_j$.
Modulo two,
\[
 \begin{split}
 M_n&\equiv n+n\sum_i p_i+B_n+P_n,\\
 E_n&\equiv P_n+(n-1)\sum_i p_i+\binom n2.
 \end{split}
\]
Since $S_n+\binom n2=(n-1)^2$, their sum gives
\[
 S_n+M_n+\delta_n\equiv\sum_i p_i+1\pmod2.
\]
This proves the common sign in
\eqref{eq:aac-current-components}.

An adjacent interchange of degrees $p,q$ changes $M_n$
by $p+q$. Reordering the corresponding parameters contributes
$(1-p)(1-q)$. Their sum is $1+pq$, which is precisely
the graded antisymmetrization sign.
Adjacent interchanges generate every permutation, so the
argument proves the formula for every arity and parity.
The cycle retains its factor $1/n$.

Define degree-zero suspended maps by
\begin{equation}\label{eq:aac-suspended-components}
 f_n(x_1\cdots x_n)
   =(-1)^{B_n(p)}s'F_n(\alpha_1,\ldots,\alpha_n)
   =(-1)^{E_n(p)}
       [t_1\cdots t_n](T_n^L+\hbar W_n^L)(A).
\end{equation}
The classical version omits $W_n$.
The signs just proved make $f_n$ graded symmetric on
$S^n(\mathcal U)$.

These maps determine a coalgebra map $\mathbf f$.
Its projection to the $k$th symmetric power of the target is
\[
 \frac1{k!}
 \sum_{\substack{i_1+\cdots+i_k=n\\i_j\geq1}}
 \ \sum_{\sigma\in\operatorname{Sh}(i_1,\ldots,i_k)}
 \epsilon_s(\sigma)\,
       f_{i_1}(x_{\sigma(1)}\cdots)
          \cdots f_{i_k}(\cdots x_{\sigma(n)}).
\]
Set $\mathbf f(1)=1$.
The scalar component vanishes on nonempty words, so the
map preserves the counit as well as the coaugmentation.
The factor $1/k!$ removes the order of the output blocks.
Splitting those blocks between two tensor factors verifies
the coalgebra identity term by term.
Every sum is finite on a finite input tensor.

\section{From the master equation to all higher identities}

Replace the variable $a$ by the assigned background weight.
In this completion, $\sum_{n\geq1}J_n$ denotes the collection
of homogeneous coefficients $J_n$ of weight $n$.
The coefficients of $d_{\mathrm{CE}}$ then equal those of
$d_{\mathrm{CE}}^{(a)}$ before this reindexing.
No evaluation of a formal series at $a=1$ occurs.

A vanishing formal curvature must imply every signed
coalgebra identity, including mixed input degrees.
Let $h_P$ denote plain graded-dual evaluation of a
coefficient-valued cochain $P$, using
\eqref{eq:aac-dual-conventions}.
Write coefficients before observables.
For a degree-zero interaction $J$ and
$v\in S^n(\mathcal U)$ of degree $r$, put $f(v)=h_J(v)$.
Then
\[
 h_{D_LJ}(v)=(-1)^rD_Lf(v),\qquad
 h_{d_{\mathrm{CE}}J}(v)=(-1)^r f(qv).
\]
For the first identity, the tensor-product differential
crosses a coefficient of degree $-r$.
For the second, the relevant coefficient before
differentiation has degree $-r-1$.
Its graded-dual sign is
$-(-1)^{-r-1}=(-1)^r$.

The bracket extended across the coefficients is
\[
 \{\phi F,\psi G\}
     =(-1)^{(|F|+1)|\psi|}\phi\psi\{F,G\}.
\]
For a split of $v$ into blocks of degrees $r_1,r_2$,
evaluation requires $|F|=r_1$ and $|\psi|=-r_2$.
The extension sign and the dual-product sign combine to
\[
 (-1)^{(r_1+1)r_2+r_1r_2}=(-1)^{r_2}.
\]
Consequently, for
$\mathcal R=(D_L+d_{\mathrm{CE}})J+\frac12\{J,J\}_L$,
\[
 h_{\mathcal R}(v)
 =(-1)^r\left(
   D_Lf(v)+f(qv)
   +\frac12\sum_{\mathrm{unshuffles}}\epsilon_s
       (-1)^{r_1}\{f(v_{(1)}),f(v_{(2)})\}_L\right).
\]
The unshuffles use two nonempty blocks.
Comparing with \eqref{eq:aac-target-q} gives the exact
corestriction identity
\begin{equation}\label{eq:aac-morphism-defect}
 \boxed{\displaystyle
 (q'\mathbf f-\mathbf f q)_{\mathrm{corestr}}(v)
          =(-1)^{r+1}h_{\mathcal R}(v).}
\end{equation}
Corestriction means projection to the first symmetric power
of the target.
Both $q'\mathbf f$ and $\mathbf f q$ are coderivations
along $\mathbf f$, as proved above.
Their difference kills the empty word.
Lemma~\ref{lem:aac-symmetric-coderivations} therefore makes
zero corestriction equivalent to the complete
coalgebra-morphism equation.

Apply this identity to the tree series, with $D_L=Q$.
Equation~\eqref{eq:aac-tree-identities} gives zero curvature.
The maps $F_n^{\mathrm{cl}}$ therefore form an
$L_\infty$ morphism from $\mathfrak g_c$ to the classical
shifted observable Lie algebra.
Here an $L_\infty$ morphism means the degree-zero coalgebra
map intertwining the specified square-zero coderivations.

In explicit components, the equation in every arity is
\begin{equation}
 \begin{split}
 q'_1f_n(x_1,\ldots,x_n)
 &+\frac12\sum_{\substack{i+j=n\\i,j\geq1}}
       \sum_{\sigma\in\operatorname{Sh}(i,j)}
       \epsilon_s(\sigma)\,
       q'_2\bigl(f_i(x_{\sigma(1)},\ldots,x_{\sigma(i)}),
                 f_j(x_{\sigma(i+1)},\ldots,x_{\sigma(n)})\bigr)\\
 &=\sum_{k=1}^n(-1)^{r_1+\cdots+r_{k-1}}
       f_n(x_1,\ldots,q_1x_k,\ldots,x_n)\\
 &\quad+\sum_{\sigma\in\operatorname{Sh}(2,n-2)}
       \epsilon_s(\sigma)\,
       f_{n-1}\bigl(q_2(x_{\sigma(1)},x_{\sigma(2)}),
                      x_{\sigma(3)},\ldots,x_{\sigma(n)}\bigr).
 \end{split}
\end{equation}
The final sum is absent for $n=1$.
Its terms respectively differentiate a source, collapse a
source edge, or join two output blocks by a heat bracket.
The general derivation \eqref{eq:aac-morphism-defect} fixes
their signs and their factor $1/2$.

At the first two arities, the quantum wheels vanish.
Writing $J(\alpha)=K_{1,L}(\alpha)$, the formulas give
\[
 \begin{split}
 F_1^q(\alpha)&=(-1)^{p+1}uJ(\alpha),\\
 F_2^q(\alpha,\eta)
   &=(-1)^{p+q+1}u\bigl(K_{2,L}(\alpha,\eta)
                            -(-1)^{pq}K_{2,L}(\eta,\alpha)\bigr).
 \end{split}
\]
These are exactly \eqref{eq:qca-noether-components}.
At arity one, $f_1(s\alpha)=(-1)^{p+1}J(\alpha)$ and
$-D_Lf_1(s\alpha)=f_1(sd\alpha)$.
At arity two, unsuspending the equation
\[
 q'_1f_2+q'_2(f_1,f_1)
   =f_2(q_1-,-)+(-1)^{p-1}f_2(-,q_1-)+f_1q_2
\]
gives
\[
 \begin{split}
 d_{\mathcal L_L}F_1(\alpha)&=F_1(d\alpha),\\
 d_{\mathcal L_L}F_2(\alpha,\eta)
   +F_2(d\alpha,\eta)+(-1)^pF_2(\alpha,d\eta)
   &=F_1([\alpha,\eta])-[F_1(\alpha),F_1(\eta)]_L.
 \end{split}
\]
No parity restriction has entered these checks.

For the unextended quantum source,
\eqref{eq:aac-all-order-qme} gives the curvature
$\hbar\Theta_\rho$. At arity three its sign can be seen
without suppressing the unshuffle terms.
For $x,y,z$ of degrees $r,s,t$, the left side is
\[
 \begin{split}
 -D_Lf_3(x,y,z)
 &+(-1)^{r+1}\{f_1x,f_2(y,z)\}_L\\
 &+(-1)^{rs+s+1}\{f_1y,f_2(x,z)\}_L\\
 &+(-1)^{t(r+s)+t+1}\{f_1z,f_2(x,y)\}_L.
 \end{split}
\]
The unextended right side is
\[
 \begin{split}
 &f_3(q_1x,y,z)+(-1)^rf_3(x,q_1y,z)
                         +(-1)^{r+s}f_3(x,y,q_1z)\\
 &\quad+f_2(q_2(x,y),z)+(-1)^{st}f_2(q_2(x,z),y)
                       +(-1)^{r(s+t)}f_2(q_2(y,z),x).
 \end{split}
\]
Their difference is the cubic value prescribed by
\eqref{eq:aac-morphism-defect}. A central source operation
will supply exactly this value.

\section{The central extension and the quantum current map}

The cubic functional is closed for the full background
differential. This consistency condition is what permits a
central extension.
Write $A_i=\partial_iA$ and
$B_0=A_1A_2-A_2A_1$.
The linear differential sends the trace integrand in
\eqref{eq:aac-theta} to $d\operatorname{Tr}(AB_0)$.
Its integral vanishes by compact support.
For the quadratic differential,
$d_{\mathrm{CE},2}A=-A^2$ and
$d_{\mathrm{CE},2}A_i=-(A_iA+AA_i)$.
The graded product rule and trace cyclicity give
\[
 d_{\mathrm{CE},2}\operatorname{Tr}(AA_1A_2)
       =\operatorname{Tr}(A^2A_1A_2).
\]
For example, expand the three differentiated factors.
The internal $AA_1AA_2$ terms cancel, and cyclically moving
the final odd $A$ gives
$\operatorname{Tr}(AA_1A_2A)=-\operatorname{Tr}(A^2A_1A_2)$.
Subtract the expression with indices exchanged.
Its integral is zero because
\[
 2\operatorname{Tr}(A^2B_0)
  =\partial_1\operatorname{Tr}(A^3A_2)
                    -\partial_2\operatorname{Tr}(A^3A_1).
\]
Expanding the derivatives cancels their mixed second
derivatives and pairs the remaining terms by graded cyclicity.
Both holomorphic divergences integrate to zero on compact
support. Therefore $d_{\mathrm{CE}}\Theta_\rho=0$.

Add a central generator $k$ of cohomological degree one.
Its suspended element $sk$ has degree zero.
Let $c$ be its degree-zero Chevalley--Eilenberg coordinate,
normalized by $c(sk)=1$.
Define the extension by
\begin{equation}\label{eq:aac-central-ce}
 d_{\mathrm{CE}}c=-\Theta_\rho(A).
\end{equation}
Its square vanishes by the cocycle identity.
There are no nonzero operations with a central input.

The actual ternary operation follows from dualization.
For homogeneous inputs put
\[
 \theta(\alpha_1,\alpha_2,\alpha_3)
       =[t_1t_2t_3]\Theta_\rho\left(\sum_i t_i\alpha_i\right).
\]
Since $|c|=0$, equation \eqref{eq:aac-dual-conventions}
gives $d_{\mathrm{CE}}c=-c\,q$.
Equations \eqref{eq:aac-central-ce} and
\eqref{eq:aac-suspended-components} therefore require
\begin{equation}\label{eq:aac-central-q}
 q_3(x_1x_2x_3)=(-1)^{E_3(p)}
            \theta(\alpha_1,\alpha_2,\alpha_3)\,sk.
\end{equation}
In unshifted graded antisymmetric notation this is
\begin{equation}
 \ell_3(\alpha_1,\alpha_2,\alpha_3)
       =(-1)^{\delta_3(p)}
            \theta(\alpha_1,\alpha_2,\alpha_3)\,k.
\end{equation}
There is no additional factorial.
The polynomial coefficient and its ordered third derivative
were distinguished in \eqref{eq:qca-polarized}.
A nonzero coefficient has $\sum_i p_i=2$, so $\ell_3$
has degree $-1$.

The extension has operations $\ell_1=d$, $\ell_2=[-,-]$,
and the displayed central $\ell_3$, with all higher
operations zero.
The corestrictions of its square at arities three and four
are the linear and quadratic cocycle identities just proved.
Those at arities one and two are unchanged.
The only possible arity-five term is a nested ternary
operation, which vanishes because its inner output is central.
There are no higher corestrictions.
The extended coderivation is odd; its square is a
coderivation with zero corestriction.
Lemma~\ref{lem:aac-symmetric-coderivations} proves that the
extended source is square-zero.

For a concrete normalization, take compact smooth functions
$f,g$, a compact smooth $(0,2)$-form $h$, and $X_0,Y_0,Z_0\in\mathfrak g$.
Put $(X_1,X_2,X_3)=(X_0,Y_0,Z_0)$ and define
\[
 T_\rho(X_0,Y_0,Z_0)
  =\frac16\sum_{\sigma\in S_3}
   \operatorname{Tr}_V
    \bigl(\rho(X_{\sigma(1)})\rho(X_{\sigma(2)})
                                    \rho(X_{\sigma(3)})\bigr).
\]
The degree pattern is $(0,0,2)$, so $E_3=\delta_3=1$
modulo two. Polarizing \eqref{eq:aac-theta} and integrating
derivatives of $h$ by parts gives
\[
 \theta(fX_0,gY_0,hZ_0)
   =-\frac1{4\pi^2}T_\rho(X_0,Y_0,Z_0)
      \int_Xh(\partial_1f\,\partial_2g
                         -\partial_2f\,\partial_1g)\wedge\Omega.
\]
The six ordered parameter terms give the factor $3!$.
Their mixed second derivatives cancel in pairs.
Thus the central operation in this degree pattern is
\[
 \ell_3(fX_0,gY_0,hZ_0)
   =\frac1{4\pi^2}T_\rho(X_0,Y_0,Z_0)
      \int_Xh(\partial_1f\,\partial_2g
                         -\partial_2f\,\partial_1g)\wedge\Omega\,k.
\]
Its positive coefficient is the consequence of the declared
central and suspension conventions.

Complete by background weight, assigning weight one to $A$
and weight three to $c$.
Write $I_L$ with its $n$th summand assigned weight $n$,
omitting the factor $a^n$.
The differential is continuous for this filtration:
its linear part preserves weight, its quadratic part
increases it, and \eqref{eq:aac-central-ce} preserves weight
three. In this completion put
\[
 \widehat I_L=I_L+\hbar c.
\]
The constant matter observable $1$ satisfies
$D_L1=0$ and $\{1,-\}_L=0$.
Equation~\eqref{eq:aac-all-order-qme} and
\eqref{eq:aac-central-ce} give
\[
 (Q+d_{\mathrm{CE}})\widehat I_L
       +\frac12\{\widehat I_L,\widehat I_L\}_L
       +\hbar\Delta_L\widehat I_L=0.
\]

\begin{theorem}\label{thm:aac-quantum-map}
At each positive heat scale $L$, the maps $F_n^q$ in
\eqref{eq:aac-current-components} extend to an algebraic
$L_\infty$ morphism from the central extension
\eqref{eq:aac-central-ce} to $\mathcal L_L$, with
\[
 F_1^q(k)=\hbar\,u1,\qquad
 F_n^q(\ldots,k,\ldots)=0\quad(n\geq2).
\]
They obey the proved bounded-family continuity estimates
on every fixed compact support sector.
\end{theorem}

\begin{proof}
Apply \eqref{eq:aac-morphism-defect} to $\widehat I_L$.
Its curvature is zero, so the coalgebra map intertwines
the two coderivations.
The linear term $\hbar c$ gives $f_1(sk)=\hbar1$ and
has no higher Taylor coefficient with a central input.
This proves the stated formulas.

The sign can also be checked against the unextended cubic
defect. A nonzero coefficient of $\Theta_\rho$ has
$r=\sum_i(p_i-1)=-1$.
Equation~\eqref{eq:aac-morphism-defect} gives
$\hbar(-1)^{E_3(p)}\theta$.
The additional source contribution $f_1q_3$ from
\eqref{eq:aac-central-q} is exactly the same value.
It cancels the defect in the extended equation.
For the degree pattern $(0,0,2)$, both values are
$-\hbar\theta$.
Every identity with a central input follows from $D_L1=0$,
centrality of $1$, and vanishing of the higher central
components.
\end{proof}

If $T_\rho=0$, then $\Theta_\rho=0$, and the quantum
maps already define a morphism from the unextended
$\mathfrak g_c$.
For an abelian representation with charges $q_j$, this
condition is $\sum_jq_j^3=0$.
For general $\rho$, removing the local central cocycle by
a local correction requires a Chevalley--Eilenberg primitive
with the specified locality and support conditions.
The construction does not assume such a primitive.

A presentation with the canonical unary suspension $-sd$
is obtained by simultaneous conjugation.
Define
\[
 U(s\alpha)=(-1)^p s\alpha,\qquad
 U'(F)=(-1)^{|F|+1}F.
\]
Conjugate both coderivations and the coalgebra map.
The degrees of their components give
\[
 \bar q_n=(-1)^nq_n,\qquad
 \bar q'_n=(-1)^nq'_n,\qquad
 \bar f_n=(-1)^{n+1}f_n.
\]
For example, $q_n$ has input degree $\sum_i(p_i-1)$
and output degree one greater. Comparing the signs of
$U$ on its output and on its $n$ inputs gives $(-1)^n$.
The same calculation gives the target and map factors.
The coordinate pullback is
$t_\alpha\mapsto(-1)^pt_\alpha$ and $c\mapsto-c$.
In particular, this presentation has
$\bar F_2=-F_2$.
Changing only the unary sign would not preserve the
displayed current identities.

The resulting action concerns compact smooth currents of
the finite-dimensional free theory on $\mathbb C^2$.
The all-arity maps require no contraction of arbitrary
distributions by $P_{0,L}$.
An identification with punctured or radial current modes
still requires the corresponding chain comparison and its
central normalization. Infinite-dimensional matter requires
its own trace and analytic estimates. Interacting vertices
change the connected graphs and require further
renormalization and master-equation arguments.
